\section*{E3. When to split the network, and where}
Two tasks share parameters $\theta$, $L(\theta)=\tfrac12\|X\theta-y_1\|^2+\tfrac12\|X\theta-y_2\|^2$, $g_i=X^\top(X\theta-y_i)$.
\begin{theorem} At a minimiser $\theta^*$ of $L$: $g_1=-g_2$, so $\cos(g_1,g_2)=-1$ unless both vanish; both vanish iff a
separate head per task loses nothing more, and the loss a split removes is
$L(\theta^*)-L_{\text{split}}=\tfrac14\|P(y_1-y_2)\|^2$, $P$ the projection on the column space of $X$.\end{theorem}
\begin{proof} $\nabla L=g_1+g_2=0$ at $\theta^*$. Write $\hat y=X\theta^*=P(y_1+y_2)/2$. Split: $\theta_i$ gives $Py_i$, loss
$\tfrac12\sum_i\|(I-P)y_i\|^2$. Shared: $\|\hat y-y_i\|^2=\|(I-P)y_i\|^2+\tfrac14\|P(y_1-y_2)\|^2$ by Pythagoras. Summing
gives the gap; $g_i=X^\top(\hat y-y_i)$ is zero for both iff $P(y_1-y_2)=0$.\end{proof}
\noindent Check. 20 seeds, three pairs: identical targets, the same direction twice the size ($y$, $2y$), opposite
($y$, $-y$). The shared loss matches the theorem to $10^{-11}$. The online rule (a moving average of the cosine,
split below $-0.9$) splits twice-size and opposite pairs in 20 of 20 runs, at step $\approx 207$ and $201$ (IQM), and
never splits identical tasks (Figure E3).

\emph{So:} a cosine stuck at $-1$ means the shared part cannot serve both. Tasks pointing the same way trip it too: they
differ only in their output. Split the \emph{head} first; the trunk only if the conflict outlives that.

\begin{figure}[h]\centering
\begin{tikzpicture}\begin{axis}[width=0.8\textwidth,height=5cm,xlabel={SGD step},ylabel={moving average of $\cos(g_1,g_2)$},ymin=-1.05,ymax=1.05,legend pos=north east]
\addplot[green!60!black,thick] table[x=step,y=identical]{../data/e3_cos.csv};
\addplot[orange,thick] table[x=step,y=twice]{../data/e3_cos.csv};
\addplot[red,thick,dashed] table[x=step,y=opposite]{../data/e3_cos.csv};
\legend{identical,twice the size,opposite}
\end{axis}\end{tikzpicture}
\caption{E3. The split signal during training (seed 0): it falls to $-1$ only where one shared output cannot fit both. Identical tasks: both gradients vanish at the optimum, the cosine is undefined and drawn $0$.}
\end{figure}
